%% maestro2tex -- generated table; do not edit by hand.
\begin{table}[htbp]
  \centering
  \caption[{Monte Carlo transimpedance loop stability, conditions as Table~\ref{tab:tiaab-mc-offset}.}]{Monte Carlo transimpedance loop stability, conditions as Table~\ref{tab:tiaab-mc-offset}. Measured by \texttt{stb} through the \texttt{IPRB\_F} probe, so the loop includes the amplifier, the feedback network and the electrode. Gain margin reports $N = 167/200$ because the loop phase never reaches the critical value on the other 33 samples and no gain margin exists for them; those samples are held out of the statistics and are not failures. Percentiles and worst case are quoted rather than \(\pm 3\sigma\) for the margins, which are not normally distributed.}
  \label{tab:tiaab-mc-stability}
  \footnotesize
  \begin{tabular}{lrrrrrrrrrr}
    \toprule
    Parameter & Unit & $N$ & Mean & $\sigma$ & $p_{1}$ & $p_{99}$ & Worst & Spec & Yield & 4\,ch \\
    \midrule
    Loop gain at DC & \si{\decibel} & 200 & 78.65 & 1.23 & 75.89 & 81.26 & 74.89 & $\geq 60.00$ & 100.0 & 100.0 \\
    Phase margin & \si{\degree} & 200 & 76.7 & 5.2 & 66.4 & 89.6 & 64.7 & $\geq 45.0$ & 100.0 & 100.0 \\
    Gain margin & \si{\decibel} & 167\,/\,200 & 22.06 & 0.64 & 20.79 & 23.70 & -- & -- & -- & -- \\
    Loop crossover & \si{\kilo\hertz} & 200 & 9.41 & 0.98 & 7.56 & 11.85 & -- & -- & -- & -- \\
    \bottomrule
  \end{tabular}
\end{table}
